% Official source: https://www.imo-official.org/assets/documents/problems/2008/2008_eng.pdf
\competitionyear{2008}

\begin{problem}{2008}{1}
Let $H$ be the orthocentre of an acute triangle $ABC$. The circle centred at the midpoint of $BC$ and passing through $H$ meets the line $BC$ at $A_1$ and $A_2$. Define $B_1, B_2$ and $C_1, C_2$ analogously using the sides $CA$ and $AB$. Prove that $A_1, A_2, B_1, B_2, C_1, C_2$ are concyclic.
\end{problem}

\begin{problem}{2008}{2}
\begin{problemsubparts}
    \item For real numbers $x, y, z$, none equal to $1$, with $xyz=1$, prove that
    \[
        \frac{x^2}{(x-1)^2}+\frac{y^2}{(y-1)^2}+\frac{z^2}{(z-1)^2}\geq1.
    \]
    \item Prove that infinitely many rational triples $(x, y, z)$ satisfying the same hypotheses attain equality.
\end{problemsubparts}
\end{problem}

\begin{problem}{2008}{3}
Prove that infinitely many positive integers $n$ have the property that $n^2+1$ possesses a prime divisor greater than $2n+\sqrt{2n}$.
\end{problem}

\begin{problem}{2008}{4}
Determine all functions $f\colon\R_{>0}\to\R_{>0}$ satisfying
\[
    \frac{(f(w))^2+(f(x))^2}{f(y^2)+f(z^2)}=\frac{w^2+x^2}{y^2+z^2}
\]
whenever $w, x, y, z$ are positive real numbers with $wx=yz$.
\end{problem}

\begin{problem}{2008}{5}
Let $n, k$ be positive integers with $k\geq n$ and $k-n$ even. There are $2n$ lamps labelled $1, \dots, 2n$, initially all off. At each step, exactly one lamp changes its state.

Let $N$ count sequences of $k$ steps ending with lamps $1, \dots, n$ on and lamps $n+1, \dots, 2n$ off. Let $M$ count those sequences with the additional restriction that lamps $n+1, \dots, 2n$ are never switched on. Determine $N/M$.
\end{problem}

\begin{problem}{2008}{6}
Let $ABCD$ be a convex quadrilateral with $BA\ne BC$. Denote the incircles of triangles $ABC$ and $ADC$ by $\omega_1$ and $\omega_2$, respectively. Suppose that a circle $\omega$ is tangent to the rays $\overrightarrow{BA}$ beyond $A$ and $\overrightarrow{BC}$ beyond $C$, as well as to the lines $AD$ and $CD$. Prove that the common external tangents to $\omega_1$ and $\omega_2$ intersect at a point on $\omega$.
\end{problem}
